\begin{lemma}[Controlled incident colors survive one iteration]
Use the zero-based translation
\(n=q+1\), \(T_n=\mathcal T.\mathrm{tree}(q+1)\),
\(\varphi_{n-1}=\mathcal T.\mathrm{coloring}(q)\), and
\(\varphi_n=\mathcal T.\mathrm{coloring}(q+1)\).  Let \(f\) be an exact
edge identity incident with \(T_n\), and suppose
\(\varphi_{n-1}(f)=c\) with
\[
 c\in\overline\varphi_{n-1}(T_n)\cup D_{n-1}.
\]
Then either \(\varphi_n(f)=\varphi_{n-1}(f)\), or the current certified
step is a PE, its singleton connecting-edge set is \(\{f\}\), its stored
\(\delta_n\) equals \(c\), and for its stored \(\gamma_n\) one has
\(\varphi_n(f)=\gamma_n\).  In all cases there is a color \(d\) such that
\[
 \varphi_n(f)=d
 \quad\hbox{and}\quad
 d\in\overline\varphi_n(T_n)\cup D_n.
\]
\end{lemma}
\begin{proof}
Eliminate the certificate at positive index \(q+1\) supplied by
\noderef{TashkinovSeries}, using the three noninitial constructors of
\noderef{TashkinovSeriesStep}.

For an RE, the constructor equality
\(\varphi_n=\varphi_{n-1}\) proves preservation immediately.  The old
color \(c\) belongs to the new missing-plus-stored set by the first
inclusion of \noderef{HkuMissingStoredMonotone}; take \(d=c\).

For an SE, write \(\pi\) for the selected coloring.  The constructor says
that \(\varphi_n=\pi\) and that \(\pi\) is
\((T_n,D_{n-1},\varphi_{n-1})\)-stable.  Since \(f\) is incident with
the tree and its old value \(c\) is in the controlling set, the first
clause of \noderef{StableColoring} gives
\(\pi(f)=\varphi_{n-1}(f)\).  Thus preservation holds, and the same
application of \noderef{HkuMissingStoredMonotone} again gives the asserted
membership with \(d=c\).

Suppose finally that the certificate is a PE.  Let \(\pi\) be the
maximum-defective selected coloring and \(\rho\) the second stable
coloring.  The first stability relation is
\((T_n,D_{n-1},\varphi_{n-1})\)-stability of \(\pi\).  The second is
\((T_n,D_{n-1}\cup\{\delta_n\},\pi)\)-stability of \(\rho\).  On
discarding the extra color \(\delta_n\), the latter remains
\((T_n,D_{n-1},\pi)\)-stable.  Both colorings are proper and leave exactly
\(e_0\) uncolored by the PE certificate.  Transitivity in
\noderef{sc2} therefore makes \(\rho\)
\((T_n,D_{n-1},\varphi_{n-1})\)-stable.  Applying its edge clause to the
given \(f,c\) yields
\[
 \rho(f)=\varphi_{n-1}(f)=c.                                    \tag{1}
\]

We next locate all edges of the switched component that are incident with
\(T_n\).  The \noderef{TreeExit} witness in the PE constructor, with the
conditions in \noderef{TreeExitPath}, gives a simple
\((\gamma_n,\delta_n)\)-path from outside \(T_n\) to its supporting
vertex \(v_n\), with no other tree vertex on the path.  The rule
\noderef{ParallelExtensionColor} says that \(\gamma_n\) is missing at
\(v_n\), whereas the exact connecting edge
\(f_n\) is incident with \(v_n\) and has color \(\delta_n\) under
\(\pi\).  Since \(\delta_n\) belongs to the enlarged controlling set,
the second stability relation fixes that edge under \(\rho\).  Hence
\(v_n\) has degree one in the two-color incidence graph and its unique
incident edge-node is \(f_n\).  The outside endpoint of the exit path has
one of the two colors missing and has the path edge of the other color, so
it also has degree one.  Every other node has degree at most two by
\noderef{TwoColorIncidenceDegreeAtMostTwo}.  If the full component had a
node outside the exit path, a shortest attachment to the path would raise
the degree of an internal path node to at least three or that of an
endpoint to at least two.  Thus the exit path is the whole component.  It
follows that \(f_n\) is the unique edge identity of this component incident
with \(T_n\).

The current coloring is the \noderef{TwoColorComponentSwap} of \(\rho\)
on precisely this component.  Therefore (1) is preserved unless
\(f=f_n\).  In that exceptional case (1), uniqueness of the incident
\(\delta_n\)-edge at \(v_n\) from \noderef{IsProperEdgeColoring}, and the
PE fields give
\[
 c=\delta_n,\qquad \varphi_n(f_n)=\gamma_n,
\]
as well as \(\Theta_n=\mathrm{PE}\) and \(F_n=\{f_n\}\).  These are
exactly the stated exception, including the identities of the edge and
both colors.

It remains to prove the unconditional membership conclusion in the PE
case.  If \(f\ne f_n\), its new color is the old color \(c\), and the
first inclusion of \noderef{HkuMissingStoredMonotone} puts \(c\) in
\(\overline\varphi_n(T_n)\cup D_n\).  If \(f=f_n\), its new color is
\(\gamma_n\), which belongs to the exact PE stored set
\(S_n=\{\delta_n,\gamma_n\}\).  Unfolding
\noderef{TashkinovStoredColors} gives
\[
 \overline\varphi_n(T_n)\cup D_n
 =\overline\varphi_n(T_n)\cup\bigcup_{h\leq n}S_h;
\]
this is the same elementary union-minus identity used for clause (i).
Hence \(\gamma_n\) belongs to the required union whether it is still
missing on the tree or lies in the set difference defining \(D_n\).
Taking \(d=\gamma_n\) completes the exceptional case and the proof.
\end{proof}
